% Generated from validated Article Draft Result. Do not hand-edit; revise the structured draft instead.
\noindent\textbf{6月のagent研究と製品更新では、複数のmodel callとtool実行をどうscheduleするか、過去の情報をどうmemoryとして保持するか、長いcontextの再処理をどう抑えるかが前面に出た。長時間agentは、model intelligenceにruntime・memory・context engineeringを重ねたsystemとして見る必要がある。}\autocite{src-80f8b4d899ff,src-c3ab33c62226,src-e8d9fd84bddc}

\subsection{Agentを一回の推論ではなく実行系として捉える}

SwarmXの著者らは、agentic workloadを独立した推論requestの集合ではなく、model callとtool executionが連鎖するworkloadとして扱い、その前提でschedulerを設計した。agentが長く動くほど、GPU計算だけでなくCPU側toolや依存関係を含むscheduleがlatencyへ効いてくるという問題設定である。\autocite{src-e8d9fd84bddc}


\subsection{永続memoryはcontext長とは別の設計問題}

OpenAIはChatGPTのmemoryを「dreaming」を用いる新architectureへ移し、reviewableなmemory summaryとstaged rolloutを組み合わせたと説明している。ここで扱われているのは単なるcontext window延長ではなく、継続利用の履歴をどの形で再構成し、次の会話へ持ち越すかというproduct-levelのstate managementである。\autocite{src-c3ab33c62226}


\subsection{context transformationを先回りする}

SmoothAgentの著者らは、長時間agentでcontext transformationが起きるたびにKV cacheが無効化され、re-prefillが発生するoverheadへ着目した。提案手法はcontext変換をahead-of-timeで進め、TTFTを抑えることを狙うsystems researchであり、報告された効果は著者のevaluation条件に依存する。\autocite{src-80f8b4d899ff}


\begin{center}
\small
\begin{tabularx}{\columnwidth}{>{\raggedright\arraybackslash}X|>{\raggedright\arraybackslash}X|>{\raggedright\arraybackslash}X}
\toprule
\textbf{層} & \textbf{6月に見えた課題} & \textbf{代表例} \\
\midrule
Scheduling & model/toolを含む複数stepの資源配分 & SwarmX \\
Persistent state & 過去情報の選別・再構成・reviewability & ChatGPT Dreaming \\
Context serving & context変換後の再処理とlatency & SmoothAgent \\
\bottomrule
\end{tabularx}
\autocite{src-80f8b4d899ff,src-c3ab33c62226,src-e8d9fd84bddc}
\end{center}

この3つをまとめると、agent性能をmodel benchmarkだけで予測することは難しい。長時間taskでは、何を覚えるか、いつtoolを動かすか、contextをどのタイミングで変換するかがsystem-levelの失敗率や応答性を左右する。model intelligenceとruntime/memory qualityは別の評価軸として持つべきである。\autocite{src-80f8b4d899ff,src-c3ab33c62226,src-e8d9fd84bddc}


\begin{claimboundary}
SwarmXとSmoothAgentの性能・latency評価は論文著者が構成した環境での報告であり、production一般で同じ改善幅を保証するものではない。問題設定とmechanismを重視し、数値は独立再現済みとして扱わない。\autocite{src-80f8b4d899ff,src-e8d9fd84bddc}
\end{claimboundary}

\begin{claimboundary}
ChatGPT Dreamingの能力やcompute efficiencyに関する説明はOpenAIのfirst-party報告である。staged rolloutやreviewable summaryという製品挙動と、効率改善の主張は分離して読む必要がある。\autocite{src-c3ab33c62226}
\end{claimboundary}
